% ============================================================
% f2013_jiangsu_q08_prism.tex
% 2013 江苏卷 第 8 题 三棱柱 A1B1C1-ABC 与截面 F-ADE
% D、E、F 分别为 AB、AC、AA1 的中点
% ============================================================
\begin{tikzpicture}[
  line width=0.6pt, font=\small, >=Stealth
]

  \coordinate (A)  at (0.2,0);
  \coordinate (B)  at (4.4,0.35);
  \coordinate (C)  at (2.0,1.95);
  \coordinate (A1) at (0.2,3.55);
  \coordinate (B1) at (4.4,3.90);
  \coordinate (C1) at (2.0,5.50);
  \coordinate (D)  at ($(A)!0.5!(B)$);
  \coordinate (E)  at ($(A)!0.5!(C)$);
  \coordinate (F)  at ($(A)!0.5!(A1)$);

  % 底面（AB、BC 可见，AC 为被遮住的边）
  \draw (A) -- (B) -- (C);
  \draw[dashed] (A) -- (C);

  % 侧棱
  \draw (A) -- (A1);
  \draw (B) -- (B1);
  \draw[dashed] (C) -- (C1);

  % 顶面
  \draw (A1) -- (B1) -- (C1) -- cycle;

  % 截面 F-ADE
  \draw (F) -- (D) -- (E);
  \draw[dashed] (E) -- (F);

  % 标注
  \node[below left]  at (A)  {$A$};
  \node[below right] at (B)  {$B$};
  \node[right]       at (C)  {$C$};
  \node[left]        at (A1) {$A_1$};
  \node[right]       at (B1) {$B_1$};
  \node[above]       at (C1) {$C_1$};
  \node[below]       at (D)  {$D$};
  \node[right]       at (E)  {$E$};
  \node[left]        at (F)  {$F$};

\end{tikzpicture}
